%% kapa: exam=1 round=24 year=2013 period=1 form=A subject=civil_law subject_ko=민법 seq=36 number=36
%% source: Q-Net 감정평가사 기출문제 (cst00309) artlSeq=5073802 file=extracted/1008/5073802/2013년 제24회 감정평가사 1차시험 기출문제/1교시(A형).hwp
%% conversion: kordoc 4.24.51 -> tools/gen1.py (auto); review=proofread-sonnet; flags=none
\begin{question}{36}
유치권에 관한 설명으로 옳은 것은? (다툼이 있으면 판례에 의함)
\begin{choices}
  \item 당사자의 합의로 유치권을 설정할 수 있다.
  \item 점유를 침탈당한 유치권자가 점유회수의 소를 제기하여 승소판결을 받았더라도 점유를 회복하지 않으면 유치권이 부활하지 않는다.
  \item 건물신축공사 수급인은 사회통념상 독립한 건물이라고 볼 수 없는 정착물에 대하여도 유치권을 취득할 수 있다.
  \item 부동산 계약명의신탁의 경우 부동산을 점유하는 명의신탁자는 명의수탁자에 대한 매매대금 상당의 부당이득반환청구권에 기해서 유치권을 행사할 수 있다.
  \item 유치권에는 물상대위가 인정된다.
\end{choices}
\end{question}
\answer{2}
